\chapter{编制目的与任务要求}\label{chap:purpose}
% REQUIREMENT-SOURCE: C:/Users/V3169/Desktop/Project/SCL_Channel_Coding_Design/任务要求/附件3-信道编码、交织与译码方案及仿真分析.pdf

\section{编制目的}

本研究面向类星导航信号体制中的低速导航电文和高速导航电文，分析电文从有效信息输入到接收端恢复之间的信道编码、译码与交织方案。研究范围及基本参数依据项目任务说明文件确定，涵盖两类业务、三类编码和七组验证任务：低速电文采用BCH码，高速电文比较卷积码与 5G NR LDPC 两类码。

不同业务对冗余、可靠性和处理时延的容忍度不同，因而需要灵活调整使用的码型。低速电文更强调接收稳健性、实现简单和译码时延可控；高速电文除可靠性外，还要考虑连续处理、有效吞吐、软判决收益以及迭代译码造成的时间波动。

多信道、译码算法和交织分别回答不同问题。多信道试验检查既定码型遇到 AWGN、多径、频偏、遮挡或突发错误时的性能变化；译码算法对比关注同一码字条件下的纠错能力、运算量、存储量和时间；交织试验只考察错误离散化对 BCH 或卷积码的改善，并记录额外缓冲和时延。

\section{业务对象与设计范围}

研究对象分为低速电文与高速电文两类。这里的“低速”和“高速”表示业务分类；

低速电文包含 \SI{200}{bit} 和 \SI{300}{bit} 两种原始有效长度，只研究 BCH。BCH码的设计同时保留分组和整块两种组织方式：分组方案把一帧电文拆成若干短 BCH 信息块，整块方案把完整信息比特适配为较长的缩短 BCH 码字。比较内容包括两种编码方式的可靠性、不同信道下的稳健性、实现复杂度和译码时间。

高速电文只保留 \SI{300}{bit}信息比特。卷积码承担连续处理和 Viterbi 译码方案，5G NR LDPC 是短块迭代译码的候选方案；两者在接近码率、相同信息比特和明确的信道条件下比较实时性、冗余和译码代价。

\section{电文输入与编码约束}

表~\ref{tab:requirement-constraints}列出输入规模、候选码型和结果形式等设计约束，并说明相应处理口径；

\begin{table}[htbp]
  \centering
  \small
  \caption{主要设计约束及报告处理方式}
  \label{tab:requirement-constraints}
  \begin{tabular}{p{0.16\textwidth}p{0.36\textwidth}p{0.38\textwidth}}
    \toprule
    项目 & 要求 & 处理方式 \\
    \midrule
    输入长度 & 低速电文约 \SI{200}{bit}、\SI{300}{bit}；高速电文只保留约 \SI{300}{bit}（第1页） & 统一按 payload 记录，附加位另列 \\
    总码长 & 单个编码块编码后总码长不超过 \SI{1000}{bit}（第1页） & 逐任务核对编码块输出和实际发送长度，不用母码参数代替 \\
    码率 & 覆盖 $1/2$、$2/3$、$3/4$、$4/5$ 附近码率（第1页） & 报告 $K_{\mathrm{payload}}/N_{\mathrm{encoded}}$ 的实际值 \\
    BCH & 分组短码与整块缩短 BCH 对比（第2--3页） & 用于 S1、S2、S6、S7 的低速候选 \\
    卷积码 & 硬/软 Viterbi，整块及按时隙处理（第3--4页） & 用于 S3、S5、S6、S7 的高速候选 \\
    LDPC & BG2 短块，BP/NMS，迭代及提前停止（第4--5页） & 用于 S4、S5、S6；Direct 实现边界单列 \\
    交织 & 独立测试；LDPC 不配置交织（第6页） & 基础曲线默认无交织，仅在 S7 突发错误场景启用 \\
    输出 & BER、FER、成功率、平均/最大时延、复杂度等曲线与表格（第7页） & 保存原始统计、参数表、图源字段及适用场景 \\
    \bottomrule
  \end{tabular}
\end{table}

码率中的“附近”意味着短块适配后允许与目标的码率存在微小的差异。BCH 缩短、卷积码尾比特和打孔、LDPC filler 或速率匹配都会改变最终长度。

\section{S1--S7仿真任务及其目的}

七项任务构成由码内筛选到专题验证的顺序。表~\ref{tab:s1-s7-purpose}列的是各任务的控制变量和输出用途，不预写尚未在对应章节审计的性能结论。

\begin{table}[htbp]
  \centering
  \scriptsize
  \caption{S1--S7 仿真任务、控制变量与用途}
  \label{tab:s1-s7-purpose}
  \begin{tabular}{p{0.06\textwidth}p{0.20\textwidth}p{0.25\textwidth}p{0.20\textwidth}p{0.20\textwidth}}
    \toprule
    任务 & 输入和候选 & 主要控制变量 & 输出 & 任务目的 \\
    \midrule
    S1 & 200/300 bit BCH，分组与整块 & 码长、码率、分组数、纠错能力 & BER、FER、成功状态、时延 & 在 BCH 内部筛选低速候选，供 S2、S6、S7 使用 \\
    S2 & S1 保留的 BCH & AWGN、多径、频偏、遮挡、突发错误 & 信道损失、FER、突发敏感性 & 检查低速候选的信道适应性 \\
    S3 & 300 bit 卷积码 & 码率、硬/软判决、量化、回溯、滑窗、时隙组织 & BER、FER、有效吞吐、首输出及译码时间 & 筛选高速卷积码候选 \\
    S4 & 300 bit BG2 Direct LDPC & 目标码长、BP/NMS、迭代上限、提前停止 & BER、FER、迭代、复杂度、时延 & 筛选高速 LDPC 候选 \\
    S5 & S3/S4 高速候选 & 接近码率、相同 payload、不同信道 & BER、FER、时延和鲁棒性 & 比较卷积码与 LDPC 的场景适用性 \\
    S6 & 固定码长与信道的三类码 & 只替换译码规则 & 性能、操作计数、存储、平均/最大时间 & 分离译码算法本身的收益和代价 \\
    S7 & 固定 BCH 或卷积码方案 & 交织器、交织深度、突发长度 & 突发容限、FER 改善、附加时延 & 评价交织；LDPC 仅作无交织基线 \\
    \bottomrule
  \end{tabular}
\end{table}

S1、S3、S4先在各自码型内部减少候选数量，避免把未经筛选的参数组合全部带入跨方案比较。S2和S5随后施加多信道条件，分别检验低速 BCH 和高速候选。S6固定编码与信道条件后只比较译码器，S7固定编码方案后只改变交织配置。这样可以把码型选择、信道适应性、译码实现和交织收益拆成可解释的四类差异。

\section{评价指标与成果形式}

\begin{table}[htbp]
  \centering
  \small
  \caption{评价指标的工程含义与使用边界}
  \label{tab:metric-categories}
  \begin{tabular}{p{0.16\textwidth}p{0.32\textwidth}p{0.42\textwidth}}
    \toprule
    类别 & 指标 & 注意 \\
    \midrule
    可靠性 & BER、FER、译码成功率、误纠、译码失败 & BER/FER 只统计恢复信息比特；\\
    传输效率 & 编码后长度、发送长度、实际码率、有效吞吐 & 附加位不进入信息比特；打孔后长度必须按实现解释 \\
    实现代价 & 操作计数、平均迭代次数、存储量、量化位宽 & 不用单一 CPU 时间替代算法复杂度 \\
    时延 & 平均、P95/P99、最大观测、首输出、结构时延、交织附加时延 & 仅在任务实际提供相应字段时报告；软件最大值不是硬件上界 \\
    场景表现 & 多信道损失、编码增益、交织改善、适用场景 & 比较点说明目标误码率、插值范围和共同控制条件 \\
    \bottomrule
  \end{tabular}
\end{table}

研究成果包括正式 BER/FER 主图、参数表、时延和复杂度表、场景说明以及可复现性记录。信噪比条件下的 BER、FER 曲线和表格以 Es/N0(符号信噪比) 为正式主图横轴。

\section{本章小结}

研究内容可归纳为四个相互关联的问题：码型和码长筛选、信道适应性、译码实现代价，以及交织对突发错误的改善。BCH、卷积码和 LDPC 分别完成码内候选筛选，再进入多信道、译码器和交织章节的比较。
